\begingroup
\fontsize{8.5pt}{9.5pt}\selectfont
\renewcommand{\arraystretch}{1.3}
\setlength{\tabcolsep}{4pt}

\begin{longtable}{|
	>{\raggedright\arraybackslash}p{1.2cm}|
	>{\raggedright\arraybackslash}p{4.0cm}|
	>{\raggedright\arraybackslash}p{2.3cm}|
	>{\raggedright\arraybackslash}p{2.0cm}|
	>{\raggedright\arraybackslash}p{4.0cm}|}

	\caption{Evaluación del cumplimiento de requisitos}
	\label{tab:evaluacion-requisitos}
	\\
	\hline
	\tableheader
	\textbf{ID}                &
	\textbf{Requisito}         &
	\textbf{Pruebas asociadas} &
	\textbf{Resultado}         &
	\textbf{Observaciones}
	\\
	\hline
	\endfirsthead

	\caption{Evaluación del cumplimiento de requisitos (continuación)}
	\\
	\hline
	\tableheader
	\textbf{ID}                &
	\textbf{Requisito}         &
	\textbf{Pruebas asociadas} &
	\textbf{Resultado}         &
	\textbf{Observaciones}
	\\
	\hline
	\endhead

	R-01                       & Gestionar centralizadamente las identidades y mantener autorización diferenciada por sistema.                                          & PT-01, PT-02                  & Cumple               & Se verificó la autenticación centralizada y la aplicación de autorización diferenciada mediante grupos.                                                                         \\ \hline
	R-02                       & Gestionar el ciclo de vida de las cuentas, incluyendo creación, actualización, habilitación, deshabilitación, revocación y expiración. & PT-03                         & Cumple               & Se verificaron las operaciones de gestión del ciclo de vida contempladas en la implementación.                                                                                  \\ \hline
	R-03                       & Definir y aplicar políticas de seguridad para usuarios y credenciales.                                                                 & PT-04                         & Cumple               & Se verificó la aplicación de las políticas de credenciales configuradas.                                                                                                        \\ \hline
	R-04                       & Controlar y monitorear los accesos y facilitar la revisión de permisos.                                                                & PT-02, PT-06                  & Cumple               & Se verificó el control de acceso mediante grupos y la revisión de los permisos asociados.                                                                                       \\ \hline
	R-05                       & Registrar actividades relevantes de gestión de identidades y accesos para fines de auditoría.                                          & PT-05                         & Cumple               & Se verificó que una operación de gestión de identidad generara un registro que permite identificar la actividad realizada.                                                      \\ \hline
	R-06                       & Fortalecer la autenticación mediante políticas de credenciales y mecanismos adicionales de autenticación.                              & PT-04                         & Cumplimiento parcial & Se implementaron políticas de credenciales, pero la integración de \acs{MFA} con los servicios mediante \acs{LDAP} no fue posible dentro de la arquitectura implementada.       \\ \hline
	R-07                       & Proporcionar mecanismos seguros de recuperación o restablecimiento de credenciales.                                                    &                               & No cumple            & FreeIPA no proporciona un mecanismo nativo de recuperación de contraseña mediante autoservicio. La incorporación de alternativas externas quedó fuera del alcance del proyecto. \\ \hline
	R-08                       & Permitir la evolución y ampliación del servicio ante la incorporación de nuevos sistemas o activos.                                    & PT-01, pruebas de integración & Cumple               & La integración de los tres servicios mediante el directorio centralizado evidencia la capacidad de incorporar sistemas al servicio de directorio.                               \\ \hline
\end{longtable}
\endgroup
\normalsize
